\section{Discussion}
\label{sec:discussion}

\paragraph{Scope, and why temperatures.}
Propositions~\ref{prop:collapse}--\ref{prop:support} say nothing about heat diffusion: they
concern matching a family of targets through a softmax readout, and apply wherever such a
family arises --- multi-teacher, multi-view or multi-layer objectives, curriculum schedules
over target sharpness. Other $r$-dependent student maps would also break the collapse
identity --- a per-scale projection head, say --- but they add parameters discarded at
inference and reintroduce the very $GL(d_S)$ freedom that Lemma~\ref{lem:gauge} identifies as
the reason pointwise anchoring fails. The temperature bank is the minimal intervention, and
by Proposition~\ref{prop:levels}(b) it already supplies calibration.

\paragraph{Limitations.}
The evidence is one teacher and one seed: Table~\ref{tab:main} has two teacher blocks
unfilled and no variance estimate, so the $+0.76$ aggregate is a not-yet-contradicted
expectation rather than a result. Every prediction in Section~\ref{sec:analysis} is stated
but untested; the SEED comparison is the most load-bearing missing cell, and the differential
structure of Prediction~\ref{pred:levels} means a merely-positive ablation would not confirm
it. $\mathcal{R}$ is bounded above by mixing rather than by choice, so whether a sparser
graph widens the usable range is open. And Condition~\ref{cond:cover} constrains the sampler
rather than the objective, so an implementation can satisfy \eqref{eq:objective} while
silently violating it.

\paragraph{Conclusion.}
The difficulty in relational distillation lies in the readout rather than in
the targets. One softmax makes a family of targets act only through its mixture, determines
the student's similarities only up to a per-anchor shift, and converts truncated targets into
a force that separates pairs the teacher calls close. A bank of per-scale temperatures over a
family that includes the zero-step, full-support target removes all three at no parameter
cost, and each claim implies a differential prediction --- what should move and what should
not --- so the experiments can falsify the analysis rather than summarise it.
